\chapter{Especificación del protocolo binario}
\label{anx:protocolo}

Formato común: \codigo{0xAA 0x55 LEN carga\_útil[LEN] CRC8}. Enteros en \emph{little-endian}.
CRC-8 con polinomio 0x07 y valor inicial 0x00 sobre la carga útil. Los desplazamientos se dan en
bytes dentro de la carga útil.

\section{Paquete 0x02: telemetría por lotes}

\begin{table}[H]
\centering
\caption{Cabecera del lote (14 bytes).}
\begin{tabularx}{\textwidth}{@{}rllX@{}}
\toprule
Despl. & Campo & Tipo & Descripción \\
\midrule
0 & packet\_id & u8 & 0x02 \\
1 & base\_t & u32 & \codigo{millis()} de la primera muestra del lote \\
5 & bridge & u16 & Lectura bruta del ZSC31014 (14 bits) \\
7 & lc\_status & u8 & Bits de estado de la última lectura \\
8 & flags & u8 & b0: célula válida; b1: servo conectado; b2: configuración aplicada \\
9 & error & u8 & Código de error (tabla~\ref{tab:errores}) \\
10 & servo\_id & u8 & Dirección Modbus en uso \\
11 & rpm\_status & u8 & Resultado de la lectura de velocidad (tabla~\ref{tab:mbstatus}) \\
12 & trq\_status & u8 & Resultado de la lectura de par \\
13 & n & u8 & Número de muestras (1--8) \\
\bottomrule
\end{tabularx}
\end{table}

\begin{table}[H]
\centering
\caption{Muestra (14 bytes), repetida $n$ veces.}
\begin{tabularx}{\textwidth}{@{}rllX@{}}
\toprule
Despl. & Campo & Tipo & Descripción \\
\midrule
+0 & dt\_ms & u16 & Milisegundos desde base\_t \\
+2 & rpm & i16 & Velocidad del motor (registro 16385) \\
+4 & current\_x10 & i16 & Par (registro 16387), décimas de \ua \\
+6 & ref\_cmd & i16 & Consigna de velocidad escrita (rpm) \\
+8 & base\_read & i16 & Lectura filtrada menos tara (cuentas) \\
+10 & io & u8 & b0: final A; b1: final B \\
+11 & state & u8 & Estado de la máquina (tabla~\ref{tab:estados}) \\
+12 & work\_us & u16 & Tiempo de trabajo de la vuelta (µs, saturado a 65535) \\
\bottomrule
\end{tabularx}
\end{table}

\section{Paquete 0x03: estado de \emph{homing} (11 bytes)}

\noindent\begin{tabularx}{\textwidth}{@{}rllX@{}}
\toprule
0 & packet\_id & u8 & 0x03 \\
1 & pos\_mrev & i32 & Posición × 1000 \\
5 & range\_mrev & i32 & Recorrido A$\rightarrow$B medido × 1000 \\
9 & home\_phase & u8 & 0 reposo, 1 buscando A, 2 buscando B, 3 yendo al origen, 4 completado, 5 fallo \\
10 & flags & u8 & b0: recorrido válido \\
\bottomrule
\end{tabularx}

\section{Paquete 0x04: resultado de la célula (13 bytes)}

\noindent\begin{tabularx}{\textwidth}{@{}rllX@{}}
\toprule
0 & packet\_id & u8 & 0x04 \\
1 & req & u8 & Orden respondida (0x09, 0x0A o 0x0B) \\
2 & result & u8 & Resultado (tabla~\ref{tab:lcres}) \\
3 & gain & u8 & Ganancia grabada (código 0--7) \\
4 & offset & u8 & Cero del A/D grabado (0--15) \\
5 & cfg\_before & u16 & B\_Config antes \\
7 & cfg\_after & u16 & B\_Config después \\
9 & ack & u8 & Primer byte devuelto por el chip \\
10 & diag & u8 & b0: el reinicio corta; b1: hubo ACK a Start\_CM; b7--b4: intento del barrido \\
11 & cfg1 & u16 & ZMDI\_Config1 \\
\bottomrule
\end{tabularx}

\section{Paquete 0x05: estado del ciclo (14 bytes)}

\noindent\begin{tabularx}{\textwidth}{@{}rllX@{}}
\toprule
0 & packet\_id & u8 & 0x05 \\
1 & phase & u8 & 0 reposo, 1 hacia B, 2 volviendo a A, 3 completado, 4 abortado \\
2 & done & u16 & Repeticiones completadas \\
4 & target & u16 & Repeticiones pedidas (0 = sin fin) \\
6 & pos\_mrev & i32 & Posición × 1000 \\
10 & target\_mrev & i32 & Posición de B × 1000 \\
\bottomrule
\end{tabularx}

\section{Órdenes}

\noindent\begin{tabularx}{\textwidth}{@{}llX@{}}
\toprule
Orden & Carga útil & Notas \\
\midrule
INIT & \codigo{01 [addr]} & addr opcional, 1--247 \\
STOP & \codigo{02} & \\
SET\_PARAM & \codigo{03 id vv vv} & valor i16 \\
MOVE\_A / MOVE\_B & \codigo{04} / \codigo{05} & \\
SHUTDOWN & \codigo{06} & \\
SCAN & \codigo{07} & bloqueante \\
HOME & \codigo{08} & \\
LC\_PROGRAM & \codigo{09 bb bb cc cc} & B\_Config u16, ZMDI\_Config1 u16 \\
LC\_QUERY & \codigo{0A} & \\
LC\_HOLD & \codigo{0B park} & 0: alta impedancia; 1: líneas a masa \\
CYCLE\_START & \codigo{0C mm tt tt vv vv rr rr} & mm = 0: B es el final; 1: B es una posición. Objetivo i16 (centi-rev), velocidad i16, repeticiones u16 \\
CYCLE\_STOP & \codigo{0D} & \\
\bottomrule
\end{tabularx}

\medskip
\noindent Atajos de terminal (sólo con el decodificador en reposo): \codigo{i} INIT, \codigo{a} MOVE\_A,
\codigo{b} MOVE\_B, \codigo{x} STOP, \codigo{h} HOME, \codigo{q} SHUTDOWN.

\section{Tablas de códigos}

\begin{table}[H]
\centering
\begin{minipage}[t]{0.47\textwidth}
\centering
\captionof{table}{Estados de la máquina.}\label{tab:estados}
\begin{tabular}{@{}rl@{}}
\toprule
0 & IDLE (servo sin habilitar) \\
1 & SCANNING \\
2 & CONFIGURING \\
3 & STOPPED (habilitado, consigna 0) \\
4 & MOVING\_A \\
5 & MOVING\_B \\
6 & ERROR \\
7 & HOMING \\
8 & PROGRAMMING (célula) \\
\bottomrule
\end{tabular}
\end{minipage}\hfill
\begin{minipage}[t]{0.47\textwidth}
\centering
\captionof{table}{Códigos de error.}\label{tab:errores}
\begin{tabular}{@{}rl@{}}
\toprule
0x00 & Sin error \\
0x01 & Célula no encontrada \\
0x02 & Fallo en el escaneo del servo \\
0x03 & Fallo al configurar el servo \\
0x04 & Límite de fuerza \\
\bottomrule
\end{tabular}

\vspace{4mm}
\captionof{table}{Resultados de ModbusMaster.}\label{tab:mbstatus}
\begin{tabular}{@{}rl@{}}
\toprule
0x00 & Correcto \\
0x01--0x04 & Excepción del esclavo \\
0xE0 & Dirección de esclavo incorrecta \\
0xE1 & Función incorrecta \\
0xE2 & Tiempo de espera agotado \\
0xE3 & CRC incorrecto \\
\bottomrule
\end{tabular}
\end{minipage}
\end{table}

\begin{table}[H]
\centering
\caption{Resultados de las órdenes de la célula.}
\label{tab:lcres}
\begin{tabular}{@{}rl@{}}
\toprule
0 & Nunca se ha intentado \\
1 & Grabado y verificado tras un ciclo de alimentación \\
2 & Ya estaba así: no se ha escrito la EEPROM \\
3 & Rechazado: el eje no está parado \\
4 & Rechazado: bus ocupado \\
5 & El chip se reinicia pero no entra en modo comando \\
6 & Escrito, pero la relectura no coincide \\
7 & Sólo lectura de la configuración \\
8 & Sigue alimentada con el reinicio activo: hay otro camino de alimentación \\
9 & La célula no contesta: sin tensión o sin bus \\
10 & Reinicio mantenido para medir \\
\bottomrule
\end{tabular}
\end{table}
